\documentclass[11pt]{article}
\usepackage[a4paper,margin=23mm]{geometry}
\usepackage{amsmath,amssymb,mathtools,bm}
\usepackage{booktabs,array,longtable,microtype,xurl,hyperref}
\hypersetup{
 colorlinks=true,linkcolor=blue,citecolor=blue,urlcolor=blue,
 pdftitle={P-54 Action, Parameter, and Convention Card},
 pdfauthor={Route-F four-dimensional mainline}
}
\setlength{\emergencystretch}{1em}

\newcommand{\Spin}{\operatorname{Spin}}
\newcommand{\tr}{\operatorname{tr}}
\newcommand{\PS}{G_{\rm PS}}
\newcommand{\SM}{G_{\rm SM}}
\newcommand{\hc}{\mathrm{h.c.}}
\newcommand{\Lag}{\mathcal L}
\newcommand{\cP}{\mathcal P}
\newcommand{\ThetaP}{\Theta_{\rm P54}}
\newcommand{\OO}{\mathcal O}

\title{P-54 Action, Parameter, and Convention Card\\
\large A fail-closed definition of the four-dimensional Route-F baseline}
\author{Route-F four-dimensional mainline}
\date{30 August 2026}

\begin{document}
\maketitle

\begin{abstract}
This card removes an ambiguity hidden in the shorthand
$54_H+126_H+10_{H,\mathbb C}$.  With only those gauge representations and no
extra symmetry, both $10_H$ and $10_H^*$ couple to the three chiral $16_F$s;
the Yukawa sector therefore contains three complex symmetric matrices.  The
predictive two-matrix model instead needs a global Peccei--Quinn symmetry and
a complex gauge singlet $S$.  We freeze that latter theory as
\texttt{P54PQ-v2}; the no-PQ theory is retained as the comparison model
\texttt{P54N}, not silently mixed with it.

The complete canonically normalized renormalizable action, tensor and charge
conventions, scalar potential, vacuum coordinates, parameter quotient, and
interfaces to the spectrum, flavor, and proton-decay calculations are given
below.  A direct charge audit finds that the often-quoted printed $\eta_1$
contraction in Eq.~(24) of Babu--Khan is incompatible with the PQ charges
printed in their Eq.~(23).  Charge conservation uniquely moves one complex
conjugation.  An independent representation-ring audit also finds the
previously omitted invariant
$\chi_7\Phi_{ij}\phi_i\phi_jS^*+\hc$.  The corrected action has three
rephasing-invariant scalar CP phases.  These are definition-level corrections: a calculation that
uses the printed monomial and simultaneously claims continuous PQ symmetry
or omits $\chi_7$ is not a calculation in \texttt{P54PQ-v2}.
\end{abstract}

\tableofcontents

\section{Status and branch identity}

The primary P-layer action is
\begin{equation}
 \boxed{\texttt{P54PQ-v2}:\quad
 \Spin(10)\times U(1)_{\rm PQ},\qquad
 3(16_F)+54_{H,\mathbb R}+126_{H,\mathbb C}
 +10_{H,\mathbb C}+1_{H,\mathbb C}.}
 \label{eq:model}
\end{equation}
The singlet is not cosmetic: after the $126_H$ expectation value, it removes
the additional spontaneously broken global direction and supplies the axion
sector.  The gauge group is $\Spin(10)$, not merely $SO(10)$, because the
matter representation is a chiral half-spinor.

Two nearby theories are explicitly outside the frozen action:
\begin{align*}
 \texttt{P54N}:&\quad 3(16_F)+54_{\mathbb R}+126_{\mathbb C}
 +10_{\mathbb C},\quad\hbox{no PQ and no singlet};\\
 \texttt{P210}:&\quad 210_H+10_H+120_H+\overline{126}_H,
 \quad\hbox{comparison/rescue only}.
\end{align*}
The first has the exact shorter field list but a different, three-matrix
Yukawa action.  The second has different vacuum and flavor freedom.  Neither
may lend parameters, thresholds, or Clebsch coefficients to
\texttt{P54PQ-v2}.

\section{Spacetime, generators, and fields}

We use Minkowski signature $\eta_{\mu\nu}=\mathrm{diag}(+,-,-,-)$,
$\epsilon^{0123}=+1$, and two-component left-handed Weyl spinors.  Long roots
have length squared two.  Hermitian generators obey
\begin{equation}
 [T^A,T^B]=i f^{ABC}T^C,\qquad
 \tr_{16}(T^AT^B)=2\delta^{AB}.
 \label{eq:generator-normalization}
\end{equation}
The gauge field and curvature are
\begin{equation}
 D_\mu=\partial_\mu-i g_{10}A_\mu^AT^A,\qquad
 F_{\mu\nu}^A=\partial_\mu A_\nu^A-\partial_\nu A_\mu^A
 +g_{10}f^{ABC}A_\mu^BA_\nu^C.
\end{equation}

Vector indices $i,j,\ldots=1,\ldots,10$ are contracted with $\delta_{ij}$.
The scalar tensors are
\begin{align}
 \Phi_{ij}&=\Phi_{ji}=\Phi_{ij}^*,& \Phi_{ii}&=0,
 &\Phi&\sim54_{\mathbb R},\label{eq:54def}\\
 \Sigma_{ijklm}&=\Sigma_{[ijklm]},& \star_{10}\Sigma&=+i\Sigma,
 &\Sigma&\sim126_{\mathbb C},\label{eq:126def}\\
 \phi_i&\in\mathbb C^{10},&&&\phi&\sim10_{\mathbb C},\\
 S&\in\mathbb C,&&&S&\sim1_{\mathbb C}.
\end{align}
Interchanging $126\leftrightarrow\overline{126}$ and the sign in
\eqref{eq:126def} everywhere is a representation-label convention, not a
second model.  Our sign is fixed by the Yukawa contraction below.

The covariant derivatives in tensor notation are
\begin{align}
 (D_\mu\Phi)_{ij}&=\partial_\mu\Phi_{ij}
 -ig_{10}A_\mu^A\big[(T^A_{10})_i{}^k\Phi_{kj}
 +(T^A_{10})_j{}^k\Phi_{ik}\big],\\
 (D_\mu\phi)_i&=\partial_\mu\phi_i
 -ig_{10}A_\mu^A(T^A_{10})_i{}^j\phi_j,\\
 (D_\mu\Sigma)_{i_1\ldots i_5}&=\partial_\mu\Sigma_{i_1\ldots i_5}
 -ig_{10}A_\mu^A\sum_{r=1}^5(T^A_{10})_{i_r}{}^j
 \Sigma_{i_1\ldots i_{r-1}j i_{r+1}\ldots i_5},\\
 D_\mu S&=\partial_\mu S.
\end{align}

The continuous PQ charges are fixed, rather than defined only up to a sign,
by
\begin{equation}
 q_{\rm PQ}(\Psi,\Phi,\Sigma,\phi,S)=(-1,0,+2,-2,-4).
 \label{eq:pqcharges}
\end{equation}
Here $\Psi_a\sim16_F$, $a=1,2,3$.  Multiplying all charges by $-1$ is of
course equivalent, but changing only one entry is not.

\section{The complete renormalizable action}

The Lagrangian is
\begin{equation}
 \Lag_{\rm P54PQ}=\Lag_{\rm gauge}+\Lag_{\rm kin}
 +\Lag_Y-V(\Phi,\Sigma,\phi,S),
 \label{eq:full-action}
\end{equation}
with
\begin{align}
 \Lag_{\rm gauge}&=-\frac14F^A_{\mu\nu}F^{A\mu\nu}
 +\frac{\theta_{10}g_{10}^2}{32\pi^2}
 F^A_{\mu\nu}\widetilde F^{A\mu\nu},\\
 \Lag_{\rm kin}&=i\Psi_a^\dagger\bar\sigma^\mu D_\mu\Psi_a
 +\frac12(D_\mu\Phi_{ij})(D^\mu\Phi_{ij})
 +\frac1{5!}(D_\mu\Sigma_{ijklm})^*(D^\mu\Sigma_{ijklm})\nonumber\\
 &\hspace{13mm}+(D_\mu\phi_i)^*(D^\mu\phi_i)
 +(\partial_\mu S)^*(\partial^\mu S),
 \label{eq:kinetic}
\end{align}
where $\widetilde F^{\mu\nu}=\epsilon^{\mu\nu\rho\sigma}F_{\rho\sigma}/2$.
Gauge fixing, Faddeev--Popov ghosts, and counterterms are introduced only
after a gauge and renormalization scheme are declared; they are not new
parameters of the gauge-invariant action.

\subsection{Yukawa sector}

Let $C_{10}$ denote the charge-conjugation matrix on the spinor indices and
$\Gamma_{i_1\ldots i_5}$ the antisymmetrized Clifford product projected onto
the self-duality selected in \eqref{eq:126def}.  Our normalization is
\begin{equation}
 -\Lag_Y=
 \frac12(Y_{10})_{ab}\,\Psi_a^TC_{10}\Gamma_i\Psi_b\,\phi_i^*
 +\frac{1}{2\,5!}(Y_{126})_{ab}\,\Psi_a^TC_{10}
 \Gamma_{ijklm}\Psi_b\,\Sigma_{ijklm}+\hc .
 \label{eq:yukawa}
\end{equation}
Both terms have total PQ charge zero.  Because
\begin{equation}
 16\otimes16=10_s\oplus120_a\oplus126_s,
\end{equation}
$Y_{10}=Y_{10}^T$ and $Y_{126}=Y_{126}^T$.  The coupling to $\phi$, as
opposed to $\phi^*$, has PQ charge $-4$ and is absent.  A $120_H$ and its
antisymmetric Yukawa matrix are absent, not set to a small fitted value.

\paragraph{Canonical-normalization correction (5 September 2026).}
An explicit Clifford-matrix contraction on the normalized P1 five-form
basis gives the raw unit-doublet coefficients $\sqrt2Y_{10}$ and
$2Y_{126}/\sqrt3$, while the singlet Majorana coefficient is
$4\sqrt2\,\sigma Y_{126}$. The previous simultaneous identifications
$h=Y_{10}$, $f=Y_{126}$ and $M_R=\sigma f$ omitted these factors.
The scalar action and radial $\sigma$ are unchanged. Define instead
\[
 h_D=\sqrt2Y_{10},\qquad f_D=\frac2{\sqrt3}Y_{126},\qquad
 f_M=4\sqrt2Y_{126}=2\sqrt6 f_D.
\]
After Pati--Salam and electroweak breaking, use $h=h_D$, $f=f_D$ in
the unit-canonical-bidoublet formulae, with vacuum coordinates
$v_{10}^{u,d},v_{126}^{u,d}$. The relative Clebsch convention is
\begin{align}
 M_u&=h v_{10}^u+f v_{126}^u,&
 M_d&=h v_{10}^d+f v_{126}^d,\label{eq:masssum1}\\
 M_D&=h v_{10}^u-3f v_{126}^u,&
 M_e&=h v_{10}^d-3f v_{126}^d,\label{eq:masssum2}\\
 M_R&=2\sqrt6\,f\,\sigma=f_M\sigma,&
 M_\nu&=M_L-M_D M_R^{-1}M_D^T.
 \label{eq:seesaw}
\end{align}
Equation \eqref{eq:seesaw} keeps both type-I and type-II terms; setting
$M_L=0$ is a later fit hypothesis, not part of this action card.
Equivalently one may use $f_M$ as the family matrix, but then each
canonical $126$ Dirac coupling contains $f_M/(2\sqrt6)$. These are
parameter translations, not a third Yukawa matrix. The CG normalization
constants are positive magnitudes; complex vertices must also carry their
scalar/spinor phase conventions. These factors and the relative $-3$
Clebsch are checked in
\path{code/verify_p54_spinor_intertwiners.py}; the global matter-conjugation
orientation and scalar/spinor relative phases must be carried consistently
when exporting the four doublet copies. No old CW numerical example may
silently mix the Dirac and Majorana conventions.

\paragraph{Phase-resolved implementation (5 September 2026).}
The preceding positive CG factors are a magnitude convention, not a
license to omit relative phases in the seesaw. The completed common
standard-matter dictionary in
\path{code/verify_p54_common_yukawa_phase.py} conjugates the scalar embedding
and spinor chirality together. It fixes the actual component phases and gives
\[
 \kappa_R=+i4\sqrt2,\qquad \kappa_{LL}=-i4\sqrt2,\qquad
 f_M^{\rm common}=i2\sqrt6 f_D,\qquad M_R=\sigma f_M^{\rm common}.
\]
Thus the $f_M$ in the positive-factor convention above and
$f_M^{\rm common}$ differ by $i$; they must not be substituted for each
other without transporting the Dirac and LL vertices as well. This $i$
is a basis convention, not a new CP parameter. In the retained four-copy
scalar order, $Y_u=c_1h_D+c_4f_D$ and
$Y_d=c_2^*h_D+c_3^*f_D$, with the usual $-3$ for the corresponding
lepton $f_D$ terms. The phase-resolved matching note records the complete
type-I/type-II covariance proof and replaces the old magnitude-only copy
map; neither note declares an experimental flavor fit.

\subsection{Scalar potential}

All repeated vector indices are summed.  With the tensor normalizations in
\eqref{eq:54def}--\eqref{eq:126def}, the full PQ-invariant potential is
\begin{equation}
 V=V_{\Phi\Sigma}+V_{\Phi\Sigma\phi}+V_S,
 \label{eq:potential-split}
\end{equation}
where
\begin{align}
V_{\Phi\Sigma}={}&-\frac{\mu^2}{2}\Phi_{ij}\Phi_{ij}
 +\frac c3\Phi_{ij}\Phi_{jk}\Phi_{ki}
 +\frac a4(\Phi_{ij}\Phi_{ij})^2
 +\frac b2\Phi_{ij}\Phi_{jk}\Phi_{kl}\Phi_{li}
 -\frac{\nu^2}{2\,5!}\Sigma_{ijklm}\Sigma^*_{ijklm}\nonumber\\
&+\frac{\lambda_0}{(2!)^2(5!)^2}
 (\Sigma_{ijklm}\Sigma^*_{ijklm})
 (\Sigma_{nopqr}\Sigma^*_{nopqr})\nonumber\\
&+\frac{\lambda_2}{(4!)^2}
 \Sigma_{ijklm}\Sigma^*_{ijkln}
 \Sigma_{opqrm}\Sigma^*_{opqrn}\nonumber\\
&+\frac{\lambda_4}{(3!)^2(2!)^2}
 \Sigma_{ijklm}\Sigma^*_{ijkno}
 \Sigma_{pqrlm}\Sigma^*_{pqrno}\nonumber\\
&+\frac{\lambda'_4}{(3!)^2}
 \Sigma_{ijklm}\Sigma^*_{ijkno}
 \Sigma_{pqrln}\Sigma^*_{pqrmo}
 +\frac{\alpha}{2!5!}\Phi_{ij}\Phi_{ij}
 \Sigma_{pqrlm}\Sigma^*_{pqrlm}\nonumber\\
&+\frac\beta{3!}\Phi_{ij}\Phi_{kl}
 \Sigma_{mnoik}\Sigma^*_{mnojl},
\label{eq:Vphisigma}
\end{align}
\begin{align}
V_{\Phi\Sigma\phi}={}&-\xi_0^2\phi_i\phi_i^*
 +\xi_1(\phi_i\phi_i^*)(\phi_j\phi_j^*)
 +\xi_2(\phi_i\phi_i)(\phi_j^*\phi_j^*)
 +\xi_3\Phi_{ij}\phi_i\phi_j^*\nonumber\\
&+\frac{\gamma_1}{4!}\Sigma_{ijklm}\Sigma^*_{ijkln}
 \phi_m\phi_n^*
 +\frac{\gamma_2}{4!}\Sigma_{ijklm}\Sigma^*_{ijkln}
 \phi_n\phi_m^*
 +\frac{\eta_0}{2}\Phi_{ij}\Phi_{ij}\phi_k\phi_k^*\nonumber\\
&+\frac{\eta_1}{(3!)^2(2!)^2}
 \underbrace{\Sigma_{ijklm}\Sigma^*_{ijkpq}
 \Sigma_{lmpqn}\phi_n}_{q_{\rm PQ}=2-2+2-2=0}
 +\frac{\eta_1^*}{(3!)^2(2!)^2}
 \Sigma^*_{ijklm}\Sigma_{ijkpq}\Sigma^*_{lmpqn}\phi_n^*\nonumber\\
&+\eta_2\Phi_{ij}\Phi_{ik}\phi_j\phi_k^*
 +\frac{\eta_3}{4!}\Sigma_{ijklm}\Sigma_{ijkln}\phi_m\phi_n
 +\frac{\eta_3^*}{4!}\Sigma^*_{ijklm}\Sigma^*_{ijkln}
 \phi_m^*\phi_n^*,
\label{eq:Vmix}
\end{align}
and
\begin{align}
V_S={}&-\mu_s^2SS^*+\chi_1(SS^*)^2
 +\chi_2\Sigma_{ijklm}\Sigma^*_{ijklm}SS^*
 +\chi_3\Phi_{ij}\Phi_{ij}SS^*\nonumber\\
&+\frac{\chi_4}{4!}\Sigma_{ijklm}\Sigma_{ijkln}\Phi_{mn}S
 +\frac{\chi_4^*}{4!}\Sigma^*_{ijklm}\Sigma^*_{ijkln}\Phi_{mn}S^*
 +\chi_5\phi_i\phi_i^*SS^*\nonumber\\
&+\chi_6\phi_i\phi_iS^*+\chi_6^*\phi_i^*\phi_i^*S
+\chi_7\Phi_{ij}\phi_i\phi_jS^*
+\chi_7^*\Phi_{ij}\phi_i^*\phi_j^*S.
\label{eq:VS}
\end{align}
The $\Phi_{mn}$ contraction in the $\chi_4$ term is forced by its two free
five-form indices.  The often-reproduced $\Phi_{ij}$ version repeats $i,j$
more than twice and is not valid Einstein notation; this card uses the
index-balanced contraction.  Couplings
$\eta_1,\eta_3,\chi_4,\chi_6,\chi_7$ are complex; all other displayed scalar
couplings are real.  Reality of the action follows from the explicit
Hermitian pairs.

The $\chi_7$ term is forced by the unique $54$ channel in
$\mathrm{Sym}^2(10)=1\oplus54$: its charge is
$q(\Phi)+2q(\phi)+q(S^*)=0-4+4=0$.  It was absent from the published-style
list used for version~1, so the independent basis audit changes the frozen
model rather than treating the omission as a numerical uncertainty.

\subsection{The corrected eta1 convention}

The printed expression in Ref.~\cite{BabuKhan2015} contains
$\Sigma\Sigma^*\Sigma^*\phi$.  With its printed charges, that monomial has
charge $-4$ and cannot occur in a continuously PQ-invariant action.  For a
quartic with three $\Sigma$-type fields and one $\phi$-type field, charge
neutrality requires
\begin{equation}
 2(n_\Sigma-n_{\Sigma^*})-2(n_\phi-n_{\phi^*})=0.
\end{equation}
The monomial inducing the $126$ bidoublet expectation value and paired with
its Hermitian conjugate is therefore uniquely
$\Sigma\Sigma^*\Sigma\phi+\hc$, which is the term in
\eqref{eq:Vmix}.  We interpret the extra star in the printed formula as a
conjugation typo.  The no-PQ branch may use the printed contraction, but then
it is not \texttt{P54PQ-v2}.

There is a second, purely index-level correction in the same printed
potential.  In $\Sigma_{ijklm}\Sigma_{ijkln}\Phi S$, the unpaired indices are
$m,n$, so the only displayed contraction is $\Phi_{mn}$, as in
\eqref{eq:VS}.  Writing $\Phi_{ij}$ would make $i,j$ occur three times and
leave $m,n$ free.  This correction changes neither PQ charge nor parameter
count.

\section{Parameter card and quotient}

Every scalar has engineering dimension one.  The coefficient card is
\begin{center}
\begin{tabular}{@{}>{\raggedright\arraybackslash}p{36mm}
 >{\raggedright\arraybackslash}p{101mm}@{\hspace{3mm}}r@{}}
\toprule
Class & Parameters & Raw reals\\
\midrule
Gauge/topological & $g_{10},\theta_{10}$ & 2\\
Scalar, dimension 2 & $\mu^2,\nu^2,\xi_0^2,\mu_s^2$ & 4\\
Scalar, dimension 1 & $c,\xi_3$ real; $\chi_6$ complex & 4\\
Scalar, dimension 0, real &
$a,b,\lambda_0,\lambda_2,\lambda_4,\lambda'_4,\alpha,\beta,
\xi_1,\xi_2,\gamma_1,\gamma_2,\eta_0,\eta_2,
\chi_1,\chi_2,\chi_3,\chi_5$ & 18\\
Scalar, dimension 0, complex & $\eta_1,\eta_3,\chi_4,\chi_7$ & 8\\
Yukawa & two complex symmetric $3\times3$ matrices & 24\\
\midrule
Total & before field-basis quotient & 60\\
\bottomrule
\end{tabular}
\end{center}

The phase-charge vectors of the five complex scalar couplings are computed
under the field redefinition
\begin{equation}
 (\Sigma,\phi,S)\longmapsto
 (e^{i\alpha_\Sigma}\Sigma,e^{i\alpha_\phi}\phi,e^{i\alpha_S}S).
\end{equation}
They are
\begin{equation}
 v_{\eta_1}=(1,1,0),\quad v_{\eta_3}=(2,2,0),\quad
 v_{\chi_4}=(2,0,1),\quad v_{\chi_6}=(0,2,-1),\quad
 v_{\chi_7}=(0,2,-1).
 \label{eq:phasevectors}
\end{equation}
Their rank is two and their common null vector is $(1,-1,-2)$, the PQ
direction in \eqref{eq:pqcharges}.  Hence only two coupling phases can be
removed.  We choose
\begin{equation}
 \eta_1\ge0,\qquad\chi_4\ge0,\qquad
 \begin{aligned}
 \delta_1&=\arg\eta_3-2\arg\eta_1,
 &\delta_2&=\arg\chi_4+\arg\chi_6-2\arg\eta_1,\\
 \delta_3&=\arg\chi_7-\arg\chi_6.&&
 \end{aligned}
 \label{eq:physicalphases}
\end{equation}
There are therefore $34-2=32$ scalar parameters after scalar-field
rephasing, including three physical scalar CP phases.

The two Yukawa matrices contain 24 real numbers.  A generic $U(3)$ family
basis removes nine; equivalently one may Takagi-diagonalize
\begin{equation}
 Y_{10}=\mathrm{diag}(y_1,y_2,y_3),\qquad y_i\ge0,
\end{equation}
and leave $Y_{126}$ complex symmetric.  This gives 15 physical Yukawa
parameters.  The classical parameter space is thus
\begin{equation}
 \dim_{\mathbb R}\bigl(\ThetaP/\mathcal G_{\rm basis}\bigr)
 =2+32+15=49.
 \label{eq:classicalcount}
\end{equation}
At the quantum level the anomalous PQ redefinition trades $\theta_{10}$ for
the axion origin, so there are 48 independent continuous parameters in
zero-temperature observables, plus discrete anomaly/domain-wall data.  This
last reduction is valid only if PQ quality is not spoiled by higher-dimension
operators.

For contrast, \texttt{P54N} has three complex symmetric Yukawa matrices,
36 raw real Yukawa entries and $36-9=27$ after the generic family quotient.
It is not a two-Yukawa-matrix realization of the shorter field list.

\section{Vacuum and breaking conventions}

The intended chain, including the global symmetry bookkeeping, is
\begin{equation}
 \Spin(10)\times U(1)_{\rm PQ}
 \xrightarrow{\langle\Phi\rangle}
 \PS\rtimes D\times U(1)_{\rm PQ}
 \xrightarrow{\langle\Sigma\rangle}
 \SM\times U(1)'_{\rm PQ}
 \xrightarrow{\langle S\rangle}\SM
 \xrightarrow{\langle h\rangle}SU(3)_C\times U(1)_{\rm em},
 \label{eq:chain}
\end{equation}
where $\PS=(SU(4)_C\times SU(2)_L\times SU(2)_R)/\mathbb Z_2$ in the
global form inherited from $\Spin(10)$.  The $54$ expectation value is
\begin{equation}
 \langle\Phi\rangle=\omega_s\,
 \mathrm{diag}\left(
 -\frac25,-\frac25,-\frac25,-\frac25,-\frac25,-\frac25,
 \frac35,\frac35,\frac35,\frac35\right).
 \label{eq:54vev}
\end{equation}
It is traceless, preserves $SO(6)\times SO(4)\simeq\PS$, and is $D$ even.
Our remaining radial coordinates are
\begin{equation}
 \langle\Sigma_{246810}\rangle=\frac{\sigma}{4\sqrt2},
 \qquad \langle S\rangle=\frac{v_s}{\sqrt2},
 \qquad \langle h^0\rangle=\frac{v}{\sqrt2},
 \label{eq:othervevs}
\end{equation}
with $\omega_s,\sigma,v_s,v\ge0$ after the stated phase conventions.

The relevant Pati--Salam decompositions are
\begin{align}
10&=(6,1,1)\oplus(1,2,2),\\
54&=(1,3,3)\oplus(6,2,2)\oplus(20',1,1)\oplus(1,1,1),\\
126&=(6,1,1)\oplus(10,3,1)\oplus(\overline{10},1,3)
 \oplus(15,2,2).
\end{align}
These labels refer to the original self-duality orientation. In the common
standard-matter dictionary of Section 3.1, the conjugated physical scalar
parents are $(\overline{10},3,1)$ for LL and $(10,1,3)$ for the
right-handed-conjugate pair; the latter carries the conjugated $\sigma$
background. No index or beta coefficient changes under this relabelling.
In the original orientation the neutral member of $(\overline{10},1,3)$
carries $\sigma$. The $54$ keeps $D$ parity, so $g_L=g_R$ at the
Pati--Salam matching scale in this branch.

Substituting \eqref{eq:54vev} into its self-potential gives the exact check
\begin{equation}
 V_{54}(\omega_s)=-\frac65\mu^2\omega_s^2
 +\frac4{25}c\omega_s^3
 +\frac{36}{25}a\omega_s^4
 +\frac{42}{125}b\omega_s^4.
 \label{eq:54radial}
\end{equation}
Indeed $\tr X^2=12/5$, $\tr X^3=12/25$, and
$\tr X^4=84/125$ for the diagonal matrix $X$ in
\eqref{eq:54vev}.  These rational coefficients are convention fingerprints:
a spectrum code with different values is not using this card.

The electroweak doublet is one normalized zero mode of the four-doublet
system
\begin{equation}
 h_{\rm SM}=c_1 H_u+c_2\epsilon H_d^*+c_3\Phi_u
 +c_4\epsilon\Phi_d^*,\qquad \sum_{i=1}^4|c_i|^2=1.
 \label{eq:lightdoublet}
\end{equation}
The $H$ doublets are from $(1,2,2)_{10}$ and the $\Phi$ doublets here are
from $(15,2,2)_{126}$; this reuse of $\Phi$ is common in Pati--Salam papers
but must not be confused with the $SO(10)$ tensor $\Phi_{ij}$ in this card.
In code they shall be named \texttt{H10u,H10d,H126u,H126d}.  The required
fine tuning is exactly one small eigenvalue of the complete doublet mass
matrix.  The color-triplet matrix must remain nonsingular and is never tuned
by hand independently.

\section{Observable map and P1 interface}

The action defines, but this card does not yet numerically evaluate, the map
\begin{equation}
 \mathcal F:\ThetaP/\mathcal G_{\rm basis}\longrightarrow\OO,
 \label{eq:observablemap}
\end{equation}
where $\OO$ contains the full scalar and vector spectrum, matched gauge
couplings, fermion singular values and mixing matrices, neutrino observables,
proton partial widths, and axion observables.  The future Jacobian is
\begin{equation}
 J_{Ai}=\frac{\partial\OO_A}{\partial\vartheta_i},\qquad
 \operatorname{rank}J\le\min(N_{\OO},48)
 \label{eq:jacobian}
\end{equation}
after the quantum PQ quotient.  No numerical rank is claimed before P1
exports a differentiable spectrum and P2--P4 fix matching conventions.
The rank must be reported from a scaled SVD with singular-value tolerance and
uncertainty propagation, not inferred from the number of fitted data.

The P1 spectrum exporter must consume exactly
\begin{equation}
 \{g_{10};\ \mu^2,\nu^2,\xi_0^2,\mu_s^2;
 c,\xi_3,\chi_6;\ \text{all dimensionless scalar couplings including }\chi_7;
 Y_{10},Y_{126}\}
\end{equation}
in the conventions above and return, for every physical state,
\begin{equation}
 \{\SM\text{ irrep},\ D\text{ parity},\ m^2,\text{mixing vector},
 \partial m^2/\partial\vartheta_i,\text{Goldstone flag}\}.
\end{equation}
Acceptance requires: all stationarity equations; equality between broken
generator directions and scalar zero modes; positive physical Hessian at the
candidate vacuum; the complete doublet and triplet matrices; and comparison
against competing stationary points.  ``Extended survival'' may be a scan
prior but may not replace this returned spectrum.

The Wilsonian action is taken at a matching scale $\mu_0$ below a declared
cutoff $\Lambda_{\rm P54}$.  The minimal consistency contract is
\begin{equation}
 M_U<\Lambda_{\rm P54}\le M_{\rm Pl},\qquad
 \max\{|\lambda_i|,|Y_{ij}|^2,g_{10}^2\}<4\pi,
 \label{eq:cutoff}
\end{equation}
supplemented in P2 by running perturbative-unitarity eigenvalues.  Planck-
suppressed PQ-violating operators are an explicit PQ-quality debt, not part
of the renormalizable fit.

\section{Freeze rules and remaining debts}

The action card is frozen under the following rules.
\begin{enumerate}
\item The corrected $\eta_1$ term in \eqref{eq:Vmix} and the $\chi_7$ term
in \eqref{eq:VS} are authoritative.  Changing either requires a new version
and a new charge audit.
\item No $120_H$, $210_H$, extra $10_H$, generic $M_R$, or arbitrary
threshold mass may be imported into the primary branch.
\item The singlet $S$ and PQ anomaly data are part of \texttt{P54PQ-v2}.
Dropping $S$ defines another theory.
\item Doublet fractions are derived eigenvectors of one scalar mass matrix;
they are not independent fit parameters in the final parameter count.
\item $M_I$, $M_U$, and every heavy mass are derived from the vacuum and
couplings.  They may be scan coordinates only with the corresponding
stationarity equations imposed.
\end{enumerate}

Version~2 is the result of the independent invariant-basis audit; its proof
and the PQ anomaly/global-form audit are recorded in the companion audit.
The machine-generated P1 Hessian/spectrum is now recorded in its own
companion and closes the tree-level algebraic vacuum gate.  The remaining
card-level debt is PQ quality above $\Lambda_{\rm P54}$; running, matching,
flavor, and proton tests remain outside this card.  The label is therefore
``tree-level candidate'', not ``phenomenologically viable theory''.

\begin{thebibliography}{9}
\bibitem{BabuKhan2015}
K.~S.~Babu and S.~Khan,
``A Minimal Non-Supersymmetric SO(10) Model: Gauge Coupling Unification,
Proton Decay and Fermion masses,''
\href{https://arxiv.org/abs/1507.06712}{arXiv:1507.06712} (2015).

\bibitem{Haba2023}
N.~Haba, Y.~Shimizu, and T.~Yamada,
``Neutrino Mass in Non-Supersymmetric SO(10) GUT,''
\href{https://arxiv.org/abs/2304.06263}{arXiv:2304.06263},
Phys. Rev. D 108, 095005 (2023).
\end{thebibliography}

\end{document}
